% Part: first-order-logic
% Chapter: beyond
% Section: introduction

\documentclass[../../../include/open-logic-section]{subfiles}

\begin{document}

\olfileid[pt]{fol}{byd}{int}

\olsection{Visão Geral}

A lógica de primeira ordem não é o único sistema de lógica de interesse:
há muitas extensões e variações da lógica de primeira ordem. Uma lógica
normalmente consiste na especificação formal de uma linguagem, geralmente,
mas nem sempre, em um sistema dedutivo e, geralmente, mas nem sempre, em
uma semântica pretendida. Mas o uso técnico do termo levanta uma
pergunta óbvia: o que as lógicas que não são lógicas de primeira ordem têm
a ver com a palavra ``lógica'', usada no sentido intuitivo ou filosófico?
Todos os sistemas descritos abaixo são projetados para modelar o
raciocínio de uma forma ou de outra; podemos dizer o que os torna lógicos?

Não há respostas fáceis. A palavra ``lógica'' é usada de
diferentes maneiras e em diferentes contextos, e a noção, como a de
``verdade'', tem sido analisada a partir de numerosos posicionamentos
filosóficos. Por exemplo, pode-se considerar que o objetivo do raciocínio
lógico é a determinação de quais afirmações são necessariamente verdadeiras,
verdadeiras a priori, verdadeiras independentemente da interpretação dos
termos não lógicos, verdadeiras em virtude de sua forma, ou verdadeiras por
convenção linguística; e cada uma dessas concepções requer bastante
esclarecimento. Mesmo que se restrinja a atenção ao tipo de lógica usada na
matemática, há pouco acordo quanto ao seu alcance. Por exemplo, no
\textit{Principia Mathematica}, Russell e Whitehead tentaram desenvolver
a matemática com base na lógica, na tradição {\em logicista}
iniciada por Frege. Seu sistema de lógica era uma forma de lógica de tipo
superior semelhante à descrita abaixo. No final, foram forçados a
introduzir axiomas que, pela maioria dos padrões, não parecem puramente
lógicos (notavelmente, o axioma do infinito e o axioma da redutibilidade),
mas pode-se, ainda assim, sustentar que algumas formas de raciocínio de
ordem superior devem ser aceitas como lógicas. Em contraste, Quine, cuja
ontologia não admite ``proposições'' como objetos legítimos de discurso,
argumenta que a lógica de segunda ordem e a de ordem superior são realmente
manifestações da teoria dos conjuntos em pele de cordeiro; em outras
palavras, sistemas que envolvem quantificação sobre predicados não são
puramente lógicos.

Por enquanto, é melhor deixar tais questões filosóficas para um dia
chuvoso e simplesmente pensar nos sistemas abaixo como idealizações formais
de vários tipos de raciocínio, lógico ou não.

\end{document}
